\subsection{希尔伯特空间的外希尔伯特和}

命题 1. --- 设 \((\mathrm{E}_i)_{i\in \mathbf{I}}\) 为一族希尔伯特空间，\(\mathbf{P}\) 为乘积向量空间 \(\prod_{i\in \mathbf{I}}\mathrm{E}_i\)，\(\mathbf{E}\) 为 \(\mathbf{P}\) 中由所有家庭 \(\mathbf{x} = (x_i)_{i\in \mathbf{I}}\)（\(\sum_{i\in \mathbf{I}}\| x_i\|^2\) 为有限）组成的子集。

\begin{enumerate}
\def\labelenumi{\alph{enumi})}
\item
  E 是 \(\mathbf{P}\) 的向量子空间。
\item
  对 E 中的每个 \(\mathbf{x} = (x_i)_{i\in \mathbf{I}}\) 与 \(\mathbf{y} = (y_i)_{i\in \mathbf{I}}\)，族 \((\langle x_i|y_i\rangle)_{i\in \mathbf{I}}\) 可和。若令 \(\langle \mathbf{x}|\mathbf{y}\rangle = \sum_{i\in \mathbf{I}}\langle x_i|y_i\rangle\)，就在 E 上定义了一个正的非退化埃尔米特型。
\item
  对于如此定义的内积，\(\mathbf{E}\) 是希尔伯特空间；\(\mathbf{S}\) 是诸 \(\mathbf{E}_i\) 的直和，它在 \(\mathbf{E}\) 中稠密。
\end{enumerate}

对 E 中的 \(\mathbf{x} = (x_{i})_{i\in \mathbf{I}}\) 与 \(\mathbf{y} = (y_i)_{i\in \mathbf{I}}\)，我们有

\[
\left\| x _ {i} + y _ {i} \right\| ^ {2} \leqslant 2 \left(\left\| x _ {i} \right\| ^ {2} + \left\| y _ {i} \right\| ^ {2}\right),
\]

因而 \(\mathbf{x} + \mathbf{y} = (x_i + y_i)_{i\in I}\) 属于 E。这就证明了 \(a)\)。

由 Cauchy-Schwarz 不等式，我们有

\[
\left| \left\langle x _ {i} \mid y _ {i} \right\rangle \right| \leqslant \| x _ {i} \| \cdot \| y _ {i} \| \leqslant \frac {1}{2} \left(\left\| x _ {i} \right\| ^ {2} + \left\| y _ {i} \right\| ^ {2}\right)
\]

因而 \(\sum_{i\in \mathbf{I}}|\langle x_i|y_i\rangle | < +\infty\)。若 \(x\neq 0\)，我们有 \(\langle \mathbf{x}|\mathbf{x}\rangle = \sum_{i\in \mathbf{I}}\| x_i\|^2 >0\)，于是断言 \(b)\) 得证。

我们回顾，S 是 P 中由所有满足如下条件的家庭 \(\mathbf{x} = (x_{i})_{i \in \mathbf{I}}\) 组成的子空间：\(i \in I\) 中使 \(x_{i} \neq 0\) 的那些所成的集是有限的。由此立即推出 S 在 E 中稠密；于是只需证明 E 对拓扑 \(T_{1}\)（由范数 \(\|x\| = \langle x|x\rangle^{1/2}\) 得到）是完备的。设 \(T_{2}\) 为 \(\prod_{i \in I} E_{i}\) 上的乘积拓扑在 E 上诱导的拓扑。对每个 r \textgreater{} 0，设 \(B_{r}\) 为 \(x \in E\) 中所有满足 \(\|x\| \leqslant r\) 的元素组成的集。此关系式蕴含：对 I 的每个有限子集 J 都有 \(\sum_{i \in J} \|x_{i}\|^{2} \leqslant r^{2}\)，因而 \(B_{r}\) 是 \(\prod_{i \in I} E_{i}\) 的闭子集，从而也是完备的。于是 E 对 \(T_{1}\) 完备这一事实可由 GT, III, §3, No.~5, 命题 10 的推论 2 得出。

定义 1. --- 设 \((\mathrm{E}_{i})_{i\in\mathrm{I}}\) 为一族希尔伯特空间。命题 1 中定义的希尔伯特空间 E 称为族 \((\mathrm{E}_{i})_{i\in\mathrm{I}}\) 的外希尔伯特和，记作 \(\bigoplus_{i\in I} E_{i}\) 或 \(\bigoplus_{i\in I} E_{i}^{1}\)。

\(^{1}\) 应注意不要把这一记号与诸空间 \(E_{i}\) 的 « 代数 » 直和的记号（A, II, § 1, No.~6）相混淆。

设 \(f_{i}\) 为从 \(E_{i}\) 到 \(E\) 中的映射，它把 \(z \in E_{i}\) 变为满足如下条件的元素 \((x_{k}) \in E\)：\(x_{k} = 0\) 对所有 \(k \neq i\) 成立，且 \(x_{i} = z\)；显然 \(f_{i}\) 是从希尔伯特空间 \(E_{i}\) 到 \(E\) 的一个闭向量子空间上的同构。我们称 \(f_{i}\) 为从 \(E_{i}\) 到 \(E\) 中的典范映射，并且通常把 \(E_{i}\) 与它在此同构下的、\(E\) 中的像等同起来。在这一约定下，\(E_{i}\) 与 \(E_{k}\) 在 \(E\) 中正交（当 \(i \neq k\) 时），而 \(E\) 是由诸子空间 \(E_{i}\) 的并生成的闭向量子空间。

当 I 为有限时，E 是诸 \(E_{i}\) 的直和；由于从 E 到 \(E_{i}\) 上的典范投影算子对所有 \(i \in I\) 都连续，E 也是诸 \(E_{i}\) 的拓扑直和（GT, III, § 6, No.~2, 命题 2）。若 \(I = \left(1, n\right)\)，我们也用 \(E_{1} \oplus E_{2} \oplus \ldots \oplus E_{n}\) 来代替 \(\bigoplus_{i=1}^{n} E_{i}\)。

例. --- 设 E 为希尔伯特空间，I 为指标集。用 \(\ell_{\mathrm{E}}^{2}(\mathrm{I})\) 表示族 \((\mathrm{E}_{i})_{i\in\mathrm{I}}\)（其中 \(E_{i}=E\) 对所有 \(i\in I\) 成立）的外希尔伯特和。换言之，\(\ell_{\mathrm{E}}^{2}(\mathrm{I})\) 是由 E 中元素的所有家庭 \(\mathbf{x}=(x_{i})_{i\in\mathrm{I}}\)（满足 \(\sum_{i\in I}\|x_{i}\|^{2}<+\infty\)）组成的空间，赋以内积 \(\langle x|y\rangle=\sum_{i\in I}\langle x_{i}|y_{i}\rangle\)（以 I 为指标集的、E 中元素之平方可和家庭所成的空间）。我们令 \(\ell^{2}(\mathrm{I})=\ell_{\mathrm{K}}^{2}(\mathrm{I})\)。

